% Generated draft: requires mathematical and source review.
\documentclass[11pt]{article}
\usepackage[margin=25mm]{geometry}
\usepackage{fontspec}
\setmainfont{Latin Modern Roman}
\usepackage{amsmath,amssymb,mathtools,unicode-math}
\setmathfont{Latin Modern Math}
\usepackage{enumitem,needspace,xurl,hyperref}
\setlength{\parindent}{0pt}
\setlength{\parskip}{5pt}
\setlength{\emergencystretch}{3em}
\newcommand{\N}{\mathbb{N}}
\newcommand{\Z}{\mathbb{Z}}
\newcommand{\Q}{\mathbb{Q}}
\newcommand{\R}{\mathbb{R}}
\newcommand{\C}{\mathbb{C}}
\newcommand{\F}{\mathbb{F}}
\newcommand{\A}{\mathbb{A}}
\newcommand{\PP}{\mathbb{P}}
\newcommand{\E}{\mathbb{E}}
\newcommand{\eps}{\varepsilon}
\newcommand{\dd}{\,\mathrm{d}}
\begin{document}
\section*{UnsolvedMath ID: 30006415}
\textbf{Source problem code:} catalogue entry
\subsection{Definitions and mathematical statement}
Let \(M\) be a smooth manifold. A Riemannian metric \(g\) on \(M\) is Einstein with Einstein constant \(\lambda\in\R\) if
\[
\operatorname{Ric}(g)=\lambda g.
\]
It is Ricci-flat when \(\lambda=0\), and positive-Einstein when \(\lambda>0\). A closed manifold is compact and without boundary. Coexistence means that two metrics are defined on the same smooth manifold; no compatibility between them is stipulated.

The supplied general question asks whether there exist an odd positive integer \(n\), a smooth closed \(n\)-manifold \(M\), Riemannian metrics \(g_0,g_1\) on \(M\), and constants \(\lambda_0,\lambda_1\in\R\) such that
\[
\operatorname{Ric}(g_i)=\lambda_i g_i\quad(i=0,1),
\qquad \lambda_0\ne\lambda_1.
\]
This wording has an unresolved normalization issue: for a constant \(a>0\), replacing an Einstein metric \(g\) by \(ag\) replaces its Einstein constant \(\lambda\) by \(\lambda/a\). Thus distinct numerical constants alone do not express a distinction invariant under constant rescaling. The supplied excerpt refers to preceding questions that are absent, and it does not specify a volume normalization, different signs, or exclusion of homothetic metrics. None of these additional conditions is imposed here.

The explicit zero-versus-positive question is precise independently of this issue: do there exist an odd positive integer \(n\), a smooth closed \(n\)-manifold \(M\), two Riemannian metrics \(g_0,g_+\) on \(M\), and a constant \(c>0\) such that
\[
\operatorname{Ric}(g_0)=0,
\qquad
\operatorname{Ric}(g_+)=c\,g_+?
\]
The record gives no lower dimension bound or connectedness assumption. It also asks what results about coexistence can be obtained using ideas related to special holonomy, without specifying a further conjecture or a required holonomy group for the general question.

As background, the supplied excerpt reports \emph{Theorem 3 (L '25)}: no smooth compact seven-manifold admits both a Sasaki--Einstein metric \(g_1\) and a metric \(g_2\) of holonomy \(G_2\). In dimension seven, a Sasaki--Einstein metric is a Sasakian Einstein metric with the usual normalization \(\operatorname{Ric}(g_1)=6g_1\); the Sasakian condition means that the metric cone \(((0,\infty)\times M,\dd r^2+r^2g_1)\) is K\"ahler. The phrase ``holonomy \(G_2\)'' is retained as equality of the Riemannian holonomy group with \(G_2\), rather than silently enlarged to holonomy contained in \(G_2\). The excerpt does not expand the attribution ``L'', clarify full versus restricted holonomy, or explicitly state a boundary convention for its use of ``compact''. This theorem is supplied background, not an independently verified assertion of this draft.
\subsection{Short English statement}
The problem asks about coexistence of Einstein metrics on the same odd-dimensional closed manifold, especially a Ricci-flat metric and a positive-Einstein metric, and whether special holonomy can provide obstructions or constructions. The broader phrase ``different Einstein constants'' needs clarification because constant rescaling changes any nonzero Einstein constant. A reported seven-dimensional Sasaki--Einstein/\(G_2\) obstruction is retained with its source qualifications.
\subsection{Sources}
\begin{itemize}
\item Supplied local UnsolvedMath registry record, permanent ID \(30006415\), in \texttt{data/unsolvedmath/problems.json}; exact source problem code \texttt{catalogue entry}. Title: \emph{Coexistence of Einstein Metrics with Different Constants}. Collection: Oberwolfach Reports Open Problems. Proposed year: 2025; no proposer is named.
\item Supplied source citation: \emph{Differentialgeometrie im Großen (2026)}, Oberwolfach Reports, DOI \texttt{10.4171/owr/2025/30}, pp.\ 1591--1646. Source URL: \url{https://doi.org/10.4171/owr/2025/30}. The supplied literature-sources entry labels the same reference \emph{Oberwolfach Report 30/2025}: \url{https://doi.org/10.4171/owr/2025/30}. The citation's 2026 title date and 2025 report designation are preserved without reconciliation.
\item The original excerpt attributes its background noncoexistence result to \emph{Theorem 3 (L '25)}. No full author name or separate theorem citation is supplied.
\item Literature assessment dated exactly \texttt{2026-08-22}: the supplied curation labels the problem open and states that the Sasaki--Einstein/\(G_2\) noncoexistence theorem for compact seven-manifolds settles one special case, while the broader odd-dimensional and zero-versus-positive questions remain open. This is supplied dated curation derived from the OWR record, not independently verified literature status or verification of the cited claims.
\item The supplied statement status is \texttt{corrected\_verified}, with a note that only formula delimiters, notation, grammar, and typography were repaired and that mathematical content and scope were preserved. This registry label does not constitute independent mathematical review of this draft.
\end{itemize}
\end{document}
